\documentclass[12pt]{article}
\usepackage[utf8]{inputenc}
\usepackage[T1]{fontenc}
\usepackage{hyperref}
\title{Tupinamb\'a Language Resources for K-12 Curriculum Development}
\author{Compiled by Claude Sonnet 5 (claude-sonnet-5) for Andr\'as Kornai}
\date{}
\begin{document}
\maketitle
\section{Introduction}
Tupinamb\'a (ISO 639-3 tpn, Glottocode tupi1273) is an extinct Tupian language of Brazil, documented in the 16th and 17th centuries by Jesuit missionaries. The descendant Nheengatu (L\'ingua Geral Amaz\^onica) is still spoken. This report is based on a single round of web searches; no dictionary or school programme was found.
\section{Primary Educational Resources}
None found. No classroom materials for Tupinamb\'a were seen in the search results.
\section{Linguistic and Reference Materials}
The main historical source is Anchieta's \emph{Arte de gram\'atica da l\'ingua mais usada na costa do Brasil} (1595). The search results discuss it, but no direct link to a full text was found. Scholarly articles that discuss it:

\url{https://periodicos.unb.br/index.php/ling/article/download/16221/14509/32031}

\url{https://revistaconfluencia.org.br/rc/article/download/1343/1126}

\url{https://seer.ufu.br/index.php/letraseletras/article/view/78165}

\url{https://www.e-publicacoes.uerj.br/soletras/article/view/87651}
\section{Religious and Community Resources}
None found.
\section{Digital and Audio Resources}
None found.
\section{Curriculum Development Context}
A heritage curriculum would rest on the 1595 grammar and other Jesuit-era texts (sources not located in this search), plus the living Nheengatu tradition. Nheengatu-based or Tupi-based revitalization programmes were not examined in this search.
\end{document}